\scanpage{88}
\mspage{75}

par
\[
\left\{\begin{aligned} W'&=\unsure{u|v}^{*-1}(W)\\ C'&=\text{clôture stable par coch. de base et composition}\\ &\hspace{1.5em}\text{transfinie de }u_!(C)
\end{aligned}\right.\tag*{(4.5.1)}
\]
\deleted{\unsure{donc}}
\[
u_!(C)\subset C'\subset\widetilde{u_!(C)}=\widetilde C'\ \deleted{\unsure{\text{ensemble}}}\tag*{\unsure{(4.5.1)|(4.5.1$'$)}}
\]
Je dis que \deleted{\unsure{$C$}} $(W',C')$, et au choix $(W',\widetilde C')$, est une structure de cofibration spéciale sur $M'$. On va vérifier les différentes conditions de 4.4. La condition a) est déjà acquise.

\margin{\uline{NB} Il vaut mieux commencer par \unsure{prouver|prendre} c)}
% slanted
b). La saturation \unsure{faible|forte} de $W'$ résulte de celle de $W$, via (4.5.1). La stabilité de $C'$ par coch. de base et par composition $\unsure{(=|C=}$ transfinie, est acquise par définition de $C'$, et idem pour $\widetilde C'$ (étant Q-close à gauche). Stabilité de $\mathrm{TC}'$\inserted{$\overset{\mathrm{déf}}{=}C'\cap W'$} par cochgt. de base ? Il suffit de voir que si on a un carré cocartésien \deleted{\unsure{de}}
\diagram{\begin{tikzcd}[ampersand replacement=\&] X'_0 \arrow[r,"i'"] \arrow[d] \arrow[dr,phantom,"\text{cocart}"] \& X'_1 \arrow[d,dash] \\ X'_2 \arrow[r,dash,"j'"] \& X'_3
\end{tikzcd}}
dans $M'$, avec $i'$ dans $\mathrm{TC}'=C'\cap W'$, alors $j'\in W'$ (car on sait déjà $j'\in C'$), i.e \margin{$\Vert$}$u^*(j')\in W\unsure{^{\prime}|{}}$. Nous allons supposer que \uline{$u^*$ commute aux sommes amalgamées}, et on est donc ramené à prouver que $u^*(i')\in\mathrm{TC}$. Donc
% archive sheet number: 87
